% Options for packages loaded elsewhere
\PassOptionsToPackage{unicode}{hyperref}
\PassOptionsToPackage{hyphens}{url}
\PassOptionsToPackage{dvipsnames,svgnames,x11names}{xcolor}
%
\documentclass[
]{article}
\usepackage{amsmath,amssymb}
\usepackage{iftex}
\ifPDFTeX
  \usepackage[T1]{fontenc}
  \usepackage[utf8]{inputenc}
  \usepackage{textcomp} % provide euro and other symbols
\else % if luatex or xetex
  \usepackage{unicode-math} % this also loads fontspec
  \defaultfontfeatures{Scale=MatchLowercase}
  \defaultfontfeatures[\rmfamily]{Ligatures=TeX,Scale=1}
\fi
\usepackage{lmodern}
\ifPDFTeX\else
  % xetex/luatex font selection
\fi
% Use upquote if available, for straight quotes in verbatim environments
\IfFileExists{upquote.sty}{\usepackage{upquote}}{}
\IfFileExists{microtype.sty}{% use microtype if available
  \usepackage[]{microtype}
  \UseMicrotypeSet[protrusion]{basicmath} % disable protrusion for tt fonts
}{}
\makeatletter
\@ifundefined{KOMAClassName}{% if non-KOMA class
  \IfFileExists{parskip.sty}{%
    \usepackage{parskip}
  }{% else
    \setlength{\parindent}{0pt}
    \setlength{\parskip}{6pt plus 2pt minus 1pt}}
}{% if KOMA class
  \KOMAoptions{parskip=half}}
\makeatother
\usepackage{xcolor}
\usepackage[margin=1in]{geometry}
\usepackage{color}
\usepackage{fancyvrb}
\newcommand{\VerbBar}{|}
\newcommand{\VERB}{\Verb[commandchars=\\\{\}]}
\DefineVerbatimEnvironment{Highlighting}{Verbatim}{commandchars=\\\{\}}
% Add ',fontsize=\small' for more characters per line
\newenvironment{Shaded}{}{}
\newcommand{\AlertTok}[1]{\textcolor[rgb]{1.00,0.00,0.00}{\textbf{#1}}}
\newcommand{\AnnotationTok}[1]{\textcolor[rgb]{0.38,0.63,0.69}{\textbf{\textit{#1}}}}
\newcommand{\AttributeTok}[1]{\textcolor[rgb]{0.49,0.56,0.16}{#1}}
\newcommand{\BaseNTok}[1]{\textcolor[rgb]{0.25,0.63,0.44}{#1}}
\newcommand{\BuiltInTok}[1]{\textcolor[rgb]{0.00,0.50,0.00}{#1}}
\newcommand{\CharTok}[1]{\textcolor[rgb]{0.25,0.44,0.63}{#1}}
\newcommand{\CommentTok}[1]{\textcolor[rgb]{0.38,0.63,0.69}{\textit{#1}}}
\newcommand{\CommentVarTok}[1]{\textcolor[rgb]{0.38,0.63,0.69}{\textbf{\textit{#1}}}}
\newcommand{\ConstantTok}[1]{\textcolor[rgb]{0.53,0.00,0.00}{#1}}
\newcommand{\ControlFlowTok}[1]{\textcolor[rgb]{0.00,0.44,0.13}{\textbf{#1}}}
\newcommand{\DataTypeTok}[1]{\textcolor[rgb]{0.56,0.13,0.00}{#1}}
\newcommand{\DecValTok}[1]{\textcolor[rgb]{0.25,0.63,0.44}{#1}}
\newcommand{\DocumentationTok}[1]{\textcolor[rgb]{0.73,0.13,0.13}{\textit{#1}}}
\newcommand{\ErrorTok}[1]{\textcolor[rgb]{1.00,0.00,0.00}{\textbf{#1}}}
\newcommand{\ExtensionTok}[1]{#1}
\newcommand{\FloatTok}[1]{\textcolor[rgb]{0.25,0.63,0.44}{#1}}
\newcommand{\FunctionTok}[1]{\textcolor[rgb]{0.02,0.16,0.49}{#1}}
\newcommand{\ImportTok}[1]{\textcolor[rgb]{0.00,0.50,0.00}{\textbf{#1}}}
\newcommand{\InformationTok}[1]{\textcolor[rgb]{0.38,0.63,0.69}{\textbf{\textit{#1}}}}
\newcommand{\KeywordTok}[1]{\textcolor[rgb]{0.00,0.44,0.13}{\textbf{#1}}}
\newcommand{\NormalTok}[1]{#1}
\newcommand{\OperatorTok}[1]{\textcolor[rgb]{0.40,0.40,0.40}{#1}}
\newcommand{\OtherTok}[1]{\textcolor[rgb]{0.00,0.44,0.13}{#1}}
\newcommand{\PreprocessorTok}[1]{\textcolor[rgb]{0.74,0.48,0.00}{#1}}
\newcommand{\RegionMarkerTok}[1]{#1}
\newcommand{\SpecialCharTok}[1]{\textcolor[rgb]{0.25,0.44,0.63}{#1}}
\newcommand{\SpecialStringTok}[1]{\textcolor[rgb]{0.73,0.40,0.53}{#1}}
\newcommand{\StringTok}[1]{\textcolor[rgb]{0.25,0.44,0.63}{#1}}
\newcommand{\VariableTok}[1]{\textcolor[rgb]{0.10,0.09,0.49}{#1}}
\newcommand{\VerbatimStringTok}[1]{\textcolor[rgb]{0.25,0.44,0.63}{#1}}
\newcommand{\WarningTok}[1]{\textcolor[rgb]{0.38,0.63,0.69}{\textbf{\textit{#1}}}}
\setlength{\emergencystretch}{3em} % prevent overfull lines
\providecommand{\tightlist}{%
  \setlength{\itemsep}{0pt}\setlength{\parskip}{0pt}}
\setcounter{secnumdepth}{-\maxdimen} % remove section numbering
\setlength{\emergencystretch}{3em}
\ifLuaTeX
  \usepackage{selnolig}  % disable illegal ligatures
\fi
\usepackage{bookmark}
\IfFileExists{xurl.sty}{\usepackage{xurl}}{} % add URL line breaks if available
\urlstyle{same}
\hypersetup{
  colorlinks=true,
  linkcolor={black},
  filecolor={Maroon},
  citecolor={Blue},
  urlcolor={black},
  pdfcreator={LaTeX via pandoc}}

\author{}
\date{}

\begin{document}

\section{DDTA research work plan after documentation-layer
closure}\label{ddta-research-work-plan-after-documentation-layer-closure}

\textbf{WORK PLAN - REVISION 6}

\textbf{Status:} active BA1 execution plan after BA1-T1 candidate
derivation.

\textbf{Supersedes:} Revision 5 only for forward execution state; R1-R5
remain historical research records.

\subsection{1. Fixed prior state}\label{1-fixed-prior-state}

\begin{itemize}
\tightlist
\item
  Chapters 2-4: CLOSED / FINAL for current thesis scope.
\item
  Documentation layer: CLOSED.
\item
  BA0-R systems-modeling prior-art research: CLOSED.
\item
  BA0 responsibility/non-goals: CLOSED by BA0-T3.
\item
  ThreatForge is a case study/tooling instrument, not DDTA semantic
  authority.
\item
  No historical focus set, systems-modeling metaclass or threat-method
  element is inherited automatically into BA1.
\end{itemize}

\subsection{2. BA0 boundary carried
forward}\label{2-ba0-boundary-carried-forward}

Base Analysis must preserve baseline-scoped shared semantic identity,
source/origin provenance and explicit unresolved state sufficiently for
reusable human/method projections, source drill-down, change-impact
reasoning and source-localized feedback handoff.

The representation remains methodology-neutral and
implementation-independent. Grounded, derived and diagnostic/unresolved
are responsibility-level origin/state distinctions; a generic reviewed
analytical-addition class is not required.

\subsection{3. BA1 objective}\label{3-ba1-objective}

BA1 decides the \textbf{minimal first-class identity ontology} needed to
realize the BA0 boundary.

For every recurring semantic responsibility, test first whether it can
remain:

\begin{itemize}
\tightlist
\item
  relation;
\item
  role/classification;
\item
  property/value;
\item
  constraint/state;
\item
  provenance/diagnostic record;
\item
  derived/projection construct.
\end{itemize}

Only independent semantic identity plus reusable subtype-specific
invariants justify a first-class split.

\begin{Shaded}
\begin{Highlighting}[]
\NormalTok{semantic responsibility exists}
\NormalTok{        !=}
\NormalTok{first{-}class metaclass required}
\end{Highlighting}
\end{Shaded}

\subsection{4. BA1-T1 - minimal ontology candidate
derivation}\label{4-ba1-t1---minimal-ontology-candidate-derivation}

\textbf{Status: COMPLETED / PROVISIONAL CANDIDATE.}

BA1-T1 rejects both extremes:

\begin{enumerate}
\def\labelenumi{\arabic{enumi}.}
\tightlist
\item
  one undifferentiated analytical element type is too weak because it
  hides the distinction between identity of a project referent and
  identity/origin/evolution of analytical claims;
\item
  domain-specific root taxonomies are premature because current evidence
  does not yet force separate \texttt{Behavior}, \texttt{Information},
  \texttt{Participant}, \texttt{Store}, \texttt{Contract},
  \texttt{Boundary}, \texttt{State}, etc. metaclasses.
\end{enumerate}

The resulting falsifiable candidate is:

\begin{Shaded}
\begin{Highlighting}[]
\NormalTok{BAReferent     CANDIDATE / NOT ACCEPTED}
\NormalTok{BAProposition  CANDIDATE / NOT ACCEPTED}
\end{Highlighting}
\end{Shaded}

\texttt{BAReferent} provides independent identity for shared project
meaning. \texttt{BAProposition} provides independent identity for
provenance-bearing shared analytical claims. Specific semantic kinds
remain open split hypotheses.

No relation vocabulary is closed by T1.

\subsection{5. BA1-T2 - split-or-collapse pressure
test}\label{5-ba1-t2---split-or-collapse-pressure-test}

\textbf{NEXT / NOT EXECUTED BY THIS PLAN REVISION.}

Pressure-test the two-family candidate against deliberately difficult
cases:

\begin{enumerate}
\def\labelenumi{\arabic{enumi}.}
\tightlist
\item
  \textbf{Behavior/event identity:} can FR-3.4 delivery and the physical
  handoff milestone be represented without a dedicated Behavior/Event
  type while preserving specialized constraints and cross-branch reuse?
\item
  \textbf{Information/resource lifecycle:} do
  \texttt{RecognitionCapture}, \texttt{OrderEvaluation},
  \texttt{Reservation}, result objects and their correlation/lifecycle
  semantics require a separate Information/Resource type?
\item
  \textbf{Responsibility and externality:} can internal/external
  ownership, provider boundaries and capability/service consumption
  remain propositions/roles over referents without dedicated
  Participant/Component/Boundary roots?
\item
  \textbf{Store and contract identity:} do the store candidate
  \texttt{InventoryLedger} and the contract candidates
  \texttt{PaymentProviderContract} / \texttt{CarrierContract} require
  subtype-specific first-class invariants?
\item
  \textbf{State/context:} can \texttt{ACCEPTED},
  \texttt{authorized/declined/indeterminate}, reservation lifecycle and
  failure states remain values/qualifiers/propositions?
\item
  \textbf{Method consumer:} can one bounded method-specific projection
  select what it needs without reconstructing missing domain kinds
  independently from raw prose?
\item
  \textbf{Collapse test:} can \texttt{BAReferent} and
  \texttt{BAProposition} be merged without losing assertion-level
  provenance/change semantics?
\end{enumerate}

\subsubsection{BA1-T2 exit}\label{ba1-t2-exit}

Record for every deferred split one of:

\begin{itemize}
\tightlist
\item
  \texttt{NOT\ FORCED};
\item
  \texttt{FORCED\ CANDIDATE\ SPLIT};
\item
  \texttt{REQUIRES\ MORE\ EVIDENCE};
\item
  \texttt{REJECTED}.
\end{itemize}

Do not close BA1 merely because the two-family candidate survives one
corpus fragment.

\subsection{6. BA1 closure review}\label{6-ba1-closure-review}

After BA1-T2, execute a distinct closure review only if the candidate is
stable enough to define a minimal ontology without embedding relation
vocabulary or provenance mechanics.

Closure must state:

\begin{itemize}
\tightlist
\item
  accepted first-class BAE types/families;
\item
  explicit non-types/roles/properties/deferred constructs;
\item
  first-class identity criterion;
\item
  reopen criteria;
\item
  what remains for BA2-BA6.
\end{itemize}

Only explicit BA1 closure authorizes BA2.

\subsection{7. Sequence after BA1
closure}\label{7-sequence-after-ba1-closure}

Do not skip phases:

\begin{Shaded}
\begin{Highlighting}[]
\NormalTok{BA2 {-} Relations and canonical action vocabulary}
\NormalTok{BA3 {-} Document{-}to{-}BAE derivation, provenance and authority}
\NormalTok{BA4 {-} Multi{-}level human and method projections}
\NormalTok{BA5 {-} Controlled vocabulary and optional authoring/extraction assistance}
\NormalTok{BA6 {-} Base Analysis closure and regression}
\end{Highlighting}
\end{Shaded}

\subsubsection{BA2}\label{ba2}

Define the smallest stable relation/action vocabulary over the accepted
BA1 identity ontology. Keep semantic relation, canonical lexical label
and authoring synonym assistance separate.

\subsubsection{BA3}\label{ba3}

Close identity/equivalence/lifecycle rules, source locators,
grounded/derived/diagnostic materialization, acceptance/rejection, stale
state and document-to-analysis traceability.

\subsubsection{BA4}\label{ba4}

Define derived human and method projections. No view is source
authority.

\subsubsection{BA5}\label{ba5}

Define controlled vocabulary and optional assistance. LLM/NLP proposals
remain candidate-producing support and cannot establish authority
silently.

\subsubsection{BA6}\label{ba6}

Regress the Base Analysis against closed corpora and at least one
structurally different holdout before formal overlay work.

\subsection{8. Analysis envelope remains
downstream}\label{8-analysis-envelope-remains-downstream}

Only after Base Analysis closure:

\begin{itemize}
\tightlist
\item
  A1 - AnalysisRecord / execution envelope;
\item
  A2 - common reviewed Finding boundary;
\item
  A3 - change/provenance integration;
\item
  formal methodology overlays and STRIDE/STRIDE-AI evaluations.
\end{itemize}

The BA0-T2 STRIDE probe remains only a pre-overlay consumer probe.

\subsection{9. Next authorized
microstep}\label{9-next-authorized-microstep}

Execute only:

\begin{quote}
\textbf{BA1-T2 - split-or-collapse pressure test of the two-family
candidate.}
\end{quote}

Do not start BA2, formal STRIDE overlay design, Common Finding schema or
implementation work until BA1 closes explicitly.

\end{document}
